\input{../tex-inputs/latex-leseansicht-vorspann}

\section[Gustav Schwarzkopf an Arthur und Olga Schnitzler, {[}14. oder 15.? 4. 1910{]}]{L04266 Gustav Schwarzkopf an Arthur und Olga Schnitzler, {[}14. oder 15.? 4. 1910{]}}
\nopagebreak\mylabel{L04266v}
\rehead{ }\normalsize\beginnumbering\briefempfaengerindex{Schnitzler, Olga@\textsc{Schnitzler, Olga}!zzzSchwarzkopf, Gustav@\emph{von Gustav Schwarzkopf}!1910-04-152@{{[}14. oder 15.? 4. 1910{]}}|(be}\briefempfaengerindex{Schnitzler, Arthur@\textsc{Schnitzler, Arthur}!zzzSchwarzkopf, Gustav@\emph{von Gustav Schwarzkopf}!1910-04-152@{{[}14. oder 15.? 4. 1910{]}}|(be}
\toendnotes[C]{\smallbreak\pagebreak[2]}
\correspDesc{Versand  durch Gustav Schwarzkopf im Zeitraum [14. oder 15.? 4. 1910] in Wien
\newline{}Erhalt  durch Arthur Schnitzler, Olga Schnitzler im Zeitraum [14. oder 15.? 4. 1910] in Wien}\toendnotes[C]{\smallbreak}
\Standort{CUL, Schnitzler, B 96a.}
\physDesc{Visitenkarte, 55 Zeichen
\newline{}Handschrift: schwarze Tinte, deutsche Kurrent
\newline{}Schnitzler: mit Bleistift datiert: »April 1910« }\toendnotes[C]{\smallbreak}
\pstart
           \centering{}{\pb}\textcolor{gray}{\textbf{Gustav Schwarzkopf}}\pend
           \vspace{0.5em}
\pstart
           gratulirt dem »\label{K_L04266-1v}\edtext{Hausherrn«
         und der »Hausfrau}{\lemma{\textnormal{\emph{Hausherrn« … »Hausfrau}}}\Cendnote{\textnormal{Schnitzler datiert
         die Visitenkarte auf »April 1910«, was gemeinsam 
            mit der Anspielung auf den am 14. 4. 1910 
            erfolgten Kaufes des Hauses Sternwartestrasse 71\oindex{Wien@\textbf{Wien}!XVIII., Währing@\textbf{XVIII., Währing}!Sternwartestraße 71@\textbf{Sternwartestraße 71}, \emph{Wohngebäude}|pwk}
            eine nähere Datierung ermöglicht. Am Folgetag am Abend trafen Schnitzler
            und Schwarzkopf\pwindex{Schwarzkopf, Gustav 7.\,11.\,1853 Wien – 13.\,11.\,1939 ebd.@\textsc{Schwarzkopf, Gustav} (7.\,11.\,1853 Wien – 13.\,11.\,1939 ebd.), \emph{Schriftsteller}|pwk} zusammen, so dass dieser die Karte
            entweder am Tag des Kaufes oder am 15. 4. 1910
            abgegeben haben dürfte.}}}\label{K_L04266-1}« herzlichſt\pend
           \selectlanguage{ngerman}\endnumbering\briefempfaengerindex{Schnitzler, Olga@\textsc{Schnitzler, Olga}!zzzSchwarzkopf, Gustav@\emph{von Gustav Schwarzkopf}!1910-04-142@{{[}14. oder 15.? 4. 1910{]}}|)be}\briefempfaengerindex{Schnitzler, Arthur@\textsc{Schnitzler, Arthur}!zzzSchwarzkopf, Gustav@\emph{von Gustav Schwarzkopf}!1910-04-142@{{[}14. oder 15.? 4. 1910{]}}|)be}\mylabel{L04266h}
\begin{anhang}
\end{anhang}\newcommand{\dateiname}{L04266}\newcommand{\titel}{Gustav Schwarzkopf an Arthur und Olga Schnitzler, [14. oder 15.? 4. 1910]}\newcommand{\editorInnen}{Herausgegeben von Jahnke, SelmaMüller, Martin Anton}\input{../tex-inputs/latex-leseansicht-abspann}
